\documentclass{article}
\usepackage{enumitem}
\usepackage{physics}
\usepackage{xcolor}

% Contadores jerárquicos (conservando tus nombres)
\newcounter{fmSubSec}[section]
\newcounter{fmCounter}[fmSubSec]
\newcommand{\theformulaCounter}{\arabic{section}.\arabic{fmSubSec}.\arabic{fmCounter}}

% Configuración GLOBAL del contador (nuevo)
\usepackage{etoolbox}
\preto\itemform{\stepcounter{fmCounter}} % Incrementa ANTES de cada ítem

% Entorno personalizado (modificado)
\newlist{itemform}{itemize}{1}
\setlist[itemform]{
  label=\textcolor{blue}{(\theformulaCounter)},
  ref=(\theformulaCounter),
  leftmargin=*,
  noitemsep,
  topsep=0pt,
  before={\setcounter{fmCounter}{0}}, % Solo reinicia al cambiar subsección
}

% Comando para subsecciones (conservando tu nombre)
\newcommand{\fmSubSec}[1]{%
  \stepcounter{fmSubSec}%
  \subsection{#1}%
}

\begin{document}

\section{Electrodinámica Clásica}

\fmSubSec{Potenciales}
\begin{itemform}
\item \label{fm:potencial-escalar} Potencial eléctrico:
\[ V(\vec{r}) = \frac{1}{4\pi\epsilon_0}\int\frac{\rho(\vec{r}')}{|\vec{r}-\vec{r}'|}d\tau' \]

\item \label{fm:campo-electrico} Campo eléctrico:
\[ \vec{E} = -\nabla V - \frac{\partial\vec{A}}{\partial t} \]
\end{itemform}

\fmSubSec{Energía}
\begin{itemform}
\item \label{fm:densidad-energia} Densidad:
\[ u = \frac{\epsilon_0}{2}E^2 + \frac{1}{2\mu_0}B^2 \]

\item \label{fm:poynting} Vector de Poynting:
\[ \vec{S} = \frac{1}{\mu_0}(\vec{E}\times\vec{B}) \]
\end{itemform}

\end{document}